\documentclass[12pt]{article}
\usepackage[T1]{fontenc}
\usepackage[utf8]{inputenc}
\usepackage[brazil]{babel}
\usepackage[margin=1in]{geometry}
\setlength{\parindent}{0pt}
\begin{document}
\section*{Tema 87}
Processo: PEDILEF 2009.72.60.000443-9/ SC\\
Situacao: Julgado (Revisado pelo Tema 128)\\
Ramo: DIREITO PREVIDENCIÁRIO\\
Questao: Saber se é possível o reconhecimento das condições especiais do labor do vigilante armado após o advento do Decreto n. 2.172/97.\\
Tese: É possível o reconhecimento de tempo especial prestado com exposição ao agente nocivo periculosidade, na atividade de vigilante, em data posterior à vigência do Decreto n. 2.172/92, de 05/03/1997, desde que laudo técnico (ou elemento material equivalente) comprove a permanente exposição à atividade nociva, com o uso de arma de fogo. Vide Tema 128 (em revisão pelo Tema 1031/STJ).\\
Fonte: https://www.cjf.jus.br/phpdoc/virtus/
\end{document}
